\honest{The Dilemma: Science or Bayes?}{Eliezer Yudkowsky}{2008}

Readers asked why I was writing so much about physics. One reader, Recovering irrationalist, wrote that many-worlds now seemed ``obvious and normal,'' and added: ``I'm just suspicious of myself for believing the first believable explanation I met.'' I was glad to see it. \nb{The part of the comment the post leaves out gives the reason: the commenter's picture of quantum physics was ``learned from one who believes overwhelmingly in MWI.'' The post does not take up that reason.}

I like having many hidden motives. It is the closest I can ethically get to being a supervillain. One of them needed physics. In healthcare or economics you may not be able to lay out formally which hypothesis is simplest or what the evidence supports; in physics the issue is clear-cut, and the evidence itself is not in dispute: ``you could write the two hypotheses out as computer programs and count the lines of code.'' \nb{The count is not given, here or in ``Decoherence is Simple,'' and what the many-worlds program must contain is disputed: a rule for probabilities, and perhaps the observer's location.} I wanted a clear example, ``Bayes says `zig', this is a zag,'' for ``when it came time to break your allegiance to Science.''

Physicists who reject many-worlds are obeying the rules of Science. The rules say a new theory must make new predictions and win a test; then come the celebration, the newspapers and the Nobel Prizes, with doddering emeritus professors quietly humoured, or else the proponent publicly recants and gains a reputation for honesty. I grant that this is how things are ``supposed'' to work, not how they do. Many-worlds does not seem to make new predictions, and its other worlds cannot be seen, which seems really suspicious, so by those rules it must wait.

A reader who started as a Traditional Rationalist and now finds many-worlds obvious should be able to flip between the two views like a Necker cube. With Science goggles on: the current theory has passed every test, many-worlds predicts nothing new we can see, it smacks of science fiction. Call me when it makes a testable prediction. With Bayes goggles: the simplest equations covering all the evidence have no exception for human-sized masses. There is no reason even to ask. Next!

Why not teach scientists Solomonoff induction? Because Science began as a rebellion against the Deep Wisdom that the wise could unravel the universe by thought alone, while looking was naive. Its core is the pragmatic belief that people in armchairs drift off into never-never land, so ``experiment alone would decide,'' not your nationality, your religion or where your hypothesis came from. Adopting many-worlds because it sounds more reasonable and elegant, without crushing the old theory in an experiment, undoes the rule that keeps people from putting angels into theories because angels seem more reasonable. \nb{That rule is one side of an old debate. Whewell defended novel predictions; Mill and Keynes rejected them. And physicists prefer special relativity to Lorentz's ether theory, which no experiment is thought able to tell apart from it, because the ether is superfluous. Preferring the simpler theory has not meant rejecting the scientific method.} Robert Aumann, who proved that Bayesian agents with similar priors cannot agree to disagree, is a believing Orthodox Jew, yet he let an experiment on the Bible codes decide, and it failed to confirm them. \nb{Accurate: Aumann wrote that research under his own supervision ``failed to confirm the existence of the codes.''}

So will you believe that parts of the wavefunction vanish, by ``the only non-linear non-unitary non-differentiable non-CPT-symmetric acausal faster-than-light informally-specified phenomenon in all of physics''? ``Or are you going to reject probability theory?'' \nb{The dilemma leaves out Bohm's theory, which has one world and no collapse, and the collapse theories that state a precise law. Probability theory does not choose the prior the simplicity argument rests on. A Bayesian can doubt many-worlds without rejecting probability theory.}

I tell Michael Vassar's joke: the arguments for majoritarianism are strong, but they must be rejected, because the majority believes in God, which is ``simply unbelievable.'' Those good at math should now be able to get ``the insanity of a global single world'' on a gut level. \nb{Two days earlier, in ``Many Worlds, One Best Guess,'' Yudkowsky had written that it is ``not quite safe'' to rank the other Earths with any other truth of science.} I wanted ``a nice, sharp dilemma between rejecting the scientific method, or embracing insanity.''

Why? ``It's not just because I'm evil''; think beyond the first obvious answer. I do not say what it is. And anyone who tries ``clever ways to wriggle out of the dilemma'' is warned that they will be shot down in future posts.
